% Done
\section{Analýza požadavků}

% Done
Na základě předchozí analýzy stávajícího řešení a analýzy konkurenčních aplikací jsem identifikoval a sepsal následující funkční a nefunkční požadavky.

% Done
\subsection{Funkční požadavky}

% Review - tohle trochu přepsat - je to napsaný moc hejhulácky
V rámci analýzy funkčních požadavků se budu zaměřovat zejména na nové funkce, které by bylo vhodné do aplikace přidat. Jelikož jsem v rámci předchozí analýzy zjistil, že aplikace obsahuje řadu funkcí, které nejsou pro uživatele tak podstatné a zbytečně komplikují celý systém, budu klást důraz na to, aby požadavky odpovídaly skutečným potřebám uživatelů a aby nezaváděly do systému zbytečnou komplexitu. Cílem je se zaměřit na užitečnost a kvalitu jednotlivých funkcí, nikoliv na jejich kvantitu.

% Write - přidat sem že Při sepisování požadavků budu také vycházet z výsledků uživatelských testů této aplikace, které jsem provedl v rámci mé bakalářské práce. 

% Review
V kapitole návrhu se pak budu věnovat podrobnému popsání případů užití, které budou jednotlivé funkce přesněji definovat.

% Review
\subsubsection{FP-1 Úvodní stránka aplikace}

% Review
\begin{tabular}{lp{0.7\linewidth}}
    \textbf{Priorita:} & Should have \\
\end{tabular}

% Review
Aplikace by měla mít úvodní stránku, která bude zobrazena, když uživatel otevře aplikaci a není přihlášen. Tato stránka by měla upoutat pozornost uživatele a umožnit mu vybrat si, jestli se chce přihlásit nebo registrovat.

% Review
Tento požadavek vychází z analýzy konkurence, kde většina aplikací má kromě stánky pro přihlášení a registraci i úvodní stránku s poutavou grafikou či textem, který prezentuje hlavní přínosy aplikace pro uživatele.

% Review
\subsubsection{FP-2 Nový úvodní dotazník a doporučování kurzů}

% Review
\begin{tabular}{lp{0.7\linewidth}}
    \textbf{Priorita:} & Should have \\
\end{tabular}

% Review
Aplikace by měla obsahovat úvodní dotazník zaměřený na získání informací o uživateli. Tyto informace by se měly týkat zejména cílů uživatele a jeho zkušeností s programováním. Na základě dotazníku by měla aplikace nabídnout uživateli několik kurzů, které odpovídají jeho cílům a zkušenostem.

% Review
Tato funkce je inspirována konkurenční aplikací Codecademy, která využívá úvodní dotazník pro doporučování kurzů uživateli.

% Review
\subsubsection{FP-3 Prohlížení a doporučování kurzů}

% Review
\begin{tabular}{lp{0.7\linewidth}}
    \textbf{Priorita:} & Must have \\
\end{tabular}

% Review
Uživatel by měl mít v aplikaci možnost zobrazit přehled dostupných kurzů a jednoduše je prohledávat. U každého kurzu by měl mít možnost se přesunout na stránku daného kurzu.

% Review
Součástí prohlížení kurzů by měla být sekce s doporučenými kurzy. Tato doporučení by měla být založena jednak na úvodním dotazníku uživatele, který vyplnil při registraci. Uživatel by měl mít možnost dotazník kdykoliv vyplnit znovu a přizpůsobit si tak svoje preference při učení. % Zde se dá inspirovat například aplikací Tinder, která má toto nastavování preferencí dobře udělané z hlediska UI/UX


% Review
\subsubsection{FP-4 Stránka kurzu}

% Review
\begin{tabular}{lp{0.7\linewidth}}
    \textbf{Priorita:} & Must have \\
\end{tabular}

% Review
Každý kurz v aplikaci by měl mít vlastní stránku s jeho popisem a přehledem lekcí, které jsou v kurzu obsaženy. Uživatel by měl mít možnost se do kurzu přihlásit nebo se z něj odhlásit. Dále by měl mít možnost u kurzu zvolit, zda se ho chce momentálně učit nebo ne. Na základě této volby by aplikace měla uživateli zobrazovat cvičení při učení.

% Review
Tento požadavek vychází z analýzy konkurence, kde ve většině aplikací mají kurzy vlastní stránku se svým popisem a přehledem lekcí.

% Review
\subsubsection{FP-5 Správa zapsaných kurzů}

% Review
\begin{tabular}{lp{0.7\linewidth}}
    \textbf{Priorita:} & Must have \\
\end{tabular}

% Review
Uživatel by měl mít možnost si zobrazit přehled zapsaných kurzů. U každého kurzu by měl uživatel vidět jeho pokrok a úroveň v rámci daného kurzu. 

% Review
Dále by u každého kurzu mělo být zobrazeno, v jakém režimu se daný kurz nachází. Režimy jsou popsány v požadavku FP-5 Učení se napříč kurzy. % Write - zde přidat odkaz na požadavek

% Review
Uživatel by měl mít také možnost se z přehledu kurzů přesunout na stránku daného kurzu.

% Review
\subsubsection{FP-6 Nastavení cílové úrovně kurzu}

% Review
\begin{tabular}{lp{0.7\linewidth}}
    \textbf{Priorita:} & Must have \\
\end{tabular}

% Review
Aplikace by měla umožnit uživateli si nastavit cílovou úroveň kurzu, tedy úroveň, které chce v uživatel daném kurzu dosáhnout. Po dosažení této úrovně by měl být kurz přepnut do režimu, v němž se uživateli bude při učení zobrazovat pouze látka k zopakování a žádné nové lekce. Uživatel si tak bude moci u každého kurzu nastavit, jaké úrovně chce dosáhnout, což mu může usnadnit učení, pokud má v aplikaci zapsaných více kurzů.

% Review
Možnost nastavení cílové úrovně by měla být k dispozici jednak při zápisu kurzu a jednak na stránce nastavení kurzu, kterou bude třeba v rámci tohoto požadavku také implementovat.

% Review
Pokud uživatel dosáhne cílové úrovně daného kurzu, mělo by se mu zobrazit oznámení, že v daném kurzu dosáhnul svého cíle a že bude kurz přepnut do režimu, v němž si uživatel bude pouze opakovat danou látku. Dále by se mu měla zobrazit informace, že režim kurzu může kdykoliv změnit v nastavení.

% Review
\subsubsection{FP-7 Kariérní cesty}

% Review
\begin{tabular}{lp{0.7\linewidth}}
    \textbf{Priorita:} & Must have \\
\end{tabular}

% Review
Aplikace by kromě kurzů měla nabízet tzv. kariérní cesty, které sdružují jednotlivé kurzy do skupin podle určitého tématu. Každá kariérní cesta reprezentuje určitou profesi. Každá kariérní cesta by měla mít vlastní stránku s popisem a přehledem kurzů, které jsou součástí dané kariérní cesty.

% Review
Uživatel by měl mít možnost si kariérní cestu zapsat a také mít možnost se od ní odhlásit.

% Review
Dále by měl mít uživatel možnosti si zvolit, zda se chce danou cestu momentálně učit nebo ne. Na základě této volby by aplikace měla uživateli zobrazovat cvičení při učení.

% Review
Tento požadavek vychází z analýzy konkurence, kde většina analyzovaných aplikací nabízí určitou formu kariérních cest, díky kterým se uživatel může učit témata týkající se konkrétní profese.

% Review
\subsubsection{FP-8 Stránka s lekcemi kurzu}

% Review
\begin{tabular}{lp{0.7\linewidth}}
    \textbf{Priorita:} & Must have \\
\end{tabular}

% Review
Ve stávající aplikaci jsou na stránce s lekcemi kurzu lekce vyobrazeny jako body na čáře připomínající cestu, po které se uživatel při postupování kurzem pohybuje. Stránka tak připomíná vizuálně mapu s cestou, po které se uživatel pohybuje. 

% Review - tyhle argumenty bych možná přehodil do analýzy stávajícího řešení
Tato mapa sice může zachytit svým vzhledem pozornost uživatele, nicméně má několik zásadních nevýhod z hlediska použitelnosti. Lekce jsou na mapě reprezentovány pouze ikonou a zkráceným názvem. Kvůli tomu může být pro uživatele obtížné na první pohled poznat, co je obsahem dané lekce. Dále, pokud se uživatel snaží najít konkrétní lekci v aplikaci, může to pro něj být velice obtížné, pokud mapa učení obsahuje větší množství lekcí.

% Review
Z těchto důvodů je třeba upravit stránku s lekcemi kurzu tak, aby byla přehlednější a umožňovala uživateli se lépe v lekcích orientovat.

% Review
\subsubsection{FP-9 Studijní plán}

% Review
\begin{tabular}{lp{0.7\linewidth}}
    \textbf{Priorita:} & Must have \\
\end{tabular}

% Review
Aplikace by měla obsahovat stránku studijního plánu uživatele. Na této stránce by měl mít uživatel možnost si zapsaných kurzů zvolit, které se chce učit a jaké úrovně v nich chce dosáhnout. Podle studijního plánu pak bude aplikace doporučovat uživateli lekce při učení.

% Review
U každého kurzu bude uživatel moci přepínat mezi následujícími režimy. Tím bude moci zvolit, které zapsané kurzy se chce učit a které ne.

% Review
\begin{itemize}
    \item \textbf{Zapnutý} -- uživateli bude aplikace zobrazovat z daného kurzu novou látku k učení i látku k zopakování
    \item \textbf{Pouze opakování} -- uživateli bude aplikace zobrazovat pouze látku k zopakování z daného kurzu
    \item \textbf{Vypnutý} -- aplikace nebude uživateli zobrazovat žádnou látku z daného kurzu
\end{itemize}

% Review
Dále bude moci u každého kurzu nastavit úroveň, které chce v kurzu dosáhnout.

% Review
Funkce studijního plánu vychází z analýzy konkurenční aplikace Codecademy, která umožňuje uživateli si vygenerovat studijní plán. Funkce dále vychází z problému, který nastává v konkurenčních aplikacích, pokud se uživatel chce učit více kurzů zároveň. Uživatel tak bývá nucen přepínat manuálně mezi jednotlivými kurzy a sám si plánovat, kterou látku se bude učit. Studijní plán by měl tomuto problému předejít a umožnit uživateli se jednoduše učit látku napříč mnoha kurzy.

\subsubsection{FP-10 Učení se podle studijního plánu}

% Review
\begin{tabular}{lp{0.7\linewidth}}
    \textbf{Priorita:} & Must have \\
\end{tabular}

% Review
Aplikace by měla uživateli poskytnout možnost se jednoduše přesunout na stránku učení. Na této stránce by měla aplikace uživateli zobrazit doporučenou lekci nebo látku k zopakování, které budou vybrány na základě znalostí uživatele a jeho studijního plánu. Lekce a látka k zopakování bude vybírána napříč všemi relevantními kurzy dle studijního plánu.

% Review
Cílem této funkce je umožnit uživateli jednoduše spustit učení z libovolného místa v aplikaci a zamezit tak problémům s přepínáním mezi kurzy, které mohou nastávat v konkurenčních aplikacích.

% Review
Tato stránka učení by zároveň měla být nastavena jako výchozí stránka, na kterou bude uživatel přesměrován po spuštění aplikace. Cílem této funkce je snížit čas od zapnutí aplikace po začátek učení a dodržet tak princip "micro-interactions", který je často využíván například v aplikacích sociálních sítí. % Write - zde dát odkaz na princip z design knížky a místo sociálnch sítí uvést příklad z té knížky + zjistit jak správně typograficky napsat micro interactions, když to je anglický pojem

% Review
\subsubsection{FP-11 Cvičení a komponenty pro výuku softwarového inženýrství} 

% Review
\begin{tabular}{lp{0.7\linewidth}}
    \textbf{Priorita:} & Must have \\
\end{tabular}

% Review
Aplikace se bude zaměřovat na výuku základů softwarového inženýrství, a proto bude obsahovat cvičení na různá témata z tohoto oboru. % Review - zde checknout, jak se píše čárka u a proto
Konkrétně by aplikace měla obsahovat následující cvičení.

% Review
\begin{itemize}
    \item Otázka s výběrem správné odpovědi
    \item Otázka a kód s vynechanými políčky, do kterých uživatel vyplní slova nebo znaky z nabídky
    \item Otázka a kód s vynechanými políčky, do kterých uživatel sám vyplní znak nebo slovo
\end{itemize}

% Review
Ve cvičeních, v nichž uživatel vyplňuje komplexnější kód, by mu měla aplikace zobrazit po vyplnění odpovědi výstup kódu. Ten může být zobrazen v následujících formátech.

% Review
\begin{itemize}
    \item Okno s jednoduchou HTML stránkou, která může obsahovat CSS styly a JavaScript kód.
    \item Okno s textem vizuálně připomínajícím výstup v terminálu.
    \item Okno s tabulkou vyplněnou daty.
\end{itemize}

% Review
Při učení bude dále možné zobrazovat uživateli následující komponenty.
\begin{itemize}
    \item Krátká ukázka kódu na jednom řádku textu
    \item Jednoduchý editor kódu s možností zobrazení více záložek reprezentující různé soubory nebo okna s výstupem kódu
    \item Obrázky ve formátech PNG, JPG, SVG a GIF
    \item Odstavce textu
    \item Nadpisy
    \item Tabulky
    \item Seznamy číslované a nečíslované
\end{itemize}

% Review
U komponent, v nichž uživatel může psát celý kód, by měla aplikace umožnit uživateli jednoduše zformátovat daný kód. Tato funkce vychází z analýzy konkurence Codecademy, která tuto funkci nabízí a může přispět k lepší uživatelské zkušenosti.

% Review
\subsubsection{FP-12 Nová struktura učení}

% Review
\begin{tabular}{lp{0.7\linewidth}}
    \textbf{Priorita:} & Must have \\
\end{tabular}

% Review
Učení by mělo být strukturováno tak, aby uživateli byla jednak po malých krocích vysvětlena probíraná látka a jednak aby si uživatel mohl danou látku procvičit. Při procvičování látky by měla aplikace sbírat data o tom, jak uživatel danému tématu rozumí. Tento způsob výuky je inspirován analyzovanými aplikacemi Mimo a Duolingo, které podobnou formu výuky využívají.

% Review - zde nakonec ověřit, jestli sedí typy komponentů 
Při učení by aplikace měla uživateli vysvětlovat látku a zadávat cvičení pomocí audio nahrávky, která bude doplněna o další komponenty, které pomohou uživateli lépe porozumět dané látce. Díky tomu bude aplikace kombinovat přístup konkurenční aplikace Mimo, která využívá různé komponenty pro vysvětlení látky, a přístup konkurenční aplikace Scrimba, v níž je látka vyučována pomocí interaktivního prostředí s audio nahrávkou vysvětlující danou látku.

% Review
Aplikace by měla brát v potaz situaci, kdy uživatel nebude moci poslouchat audio a měla by uživateli umožnit se učit i bez něj.

% Review
Tato struktura učení se zásadně liší od struktury učení v původní aplikaci, kde probíraná látka nebyla uživateli při učení vysvětlena a učení nebylo doplněno o audio nahrávku vysvětlující danou látku.

% Review
Při řešení tohoto požadavku je třeba také brát v potaz výsledky uživatelského testování stávajícího řešení, které bylo provedené v rámci mé bakalářské práce. Výsledky tohoto testování ukázaly, že aplikace nedává uživateli dostatečně jasnou zpětnou vazbu o tom, zda-li odpověděl při plnění cvičení správně nebo špatně. To bude třeba zohlednit při návrhu a implementaci nového způsobu učení.

% Review
\subsubsection{FP-13 Chatbot}

% Review
\begin{tabular}{lp{0.7\linewidth}}
    \textbf{Priorita:} & Should have \\
\end{tabular}

% Review
Aplikace by měla při učení poskytovat premium uživateli chatbota, který mu bude moci vysvětlit, proč byla jeho odpověď ve cvičení špatná nebo mu podrobněji vysvětlit, co daný kód ve cvičení znamená a jak funguje.

% Review
\subsubsection{FP-14 Nahlašování problému při učení}

% Review
\begin{tabular}{lp{0.7\linewidth}}
    \textbf{Priorita:} & Should have \\
\end{tabular}

% Review
Uživatel by měl mít při učení možnost u každého cvičení nebo vysvětlení látky nahlásit problém, který při dané aktivitě nastal. Při nahlašování by měl mít možnost zvolit o který typ problému z následujícího výběru se jedná.

% Review
\begin{itemize}
    \item Uživatel považuje svou odpověď ve cvičení za správnou, i přes to, že ji aplikace označila jako špatnou.
    \item Uživatel považuje svou odpověď ve cvičení za špatnou, i přes to, že ji aplikace označila jako správnou.
    \item Zadání cvičení je pro uživatele matoucí nebo nejasné.
    \item Vysvětlení látky je pro uživatele matoucí nebo nejasné.
    \item Nastala jiná chyba. V tomto případě by měl uživatel mít možnost sám uvést, o jakou chybu se jedná.
\end{itemize}

% Review
Uživatelé s rolí korektora by měli kromě předchozích možností nahlášení problému, mít také k dispozici následující možnosti.

% Review
\begin{itemize}
    \item Vysvětlení látky je příliš dlouhé nebo příliš krátké.
    \item Cvičení je příliš těžké nebo příliš snadné.
    \item Cvičení nedostatečně testuje probíraný koncept.
    \item Struktura lekce by měla být upravena. Tato možnost se využije například pokud v lekci není dostatek cvičení nebo vysvětlení látky.
\end{itemize}

% Review
Korektoři by také měli mít možnost u každé odpovědi detailněji popsat vlastními slovy daný problém. 

% Review
Nahlášení problému od uživatelů s rolí korektora by mělo být označeno tak, aby bylo možné jej rozlišit od nahlášení problému běžných uživatelů. Role korektora bude nastavována uživatelům mimo frontend této aplikace. Tuto roli využijí zejména lidé zodpovědní za tvorbu a korekci jednotlivých kurzů.

% Review
\subsubsection{FP-15 Statistiky o učení}

% Review
\begin{tabular}{lp{0.7\linewidth}}
    \textbf{Priorita:} & Could have \\
\end{tabular}

% Review
Aplikace by měla při učení uživatele nově sbírat data o tom, jak dlouho uživatel prováděl jednotlivá cvičení nebo procházel vysvětlení látky. Tato data by měla být využita k tomu, aby se na konci učení uživateli zobrazila informace, jak dlouho se učil. 

% Review
Dále by měla být tato data odeslána do backendové části aplikace, kde budou následně zpracována a použita pro úpravu odhadů časové náročnosti jednotlivých cvičení. Touto funkcí backendové části aplikace se zabývá kolega Bc. Jakub Vondráček v jeho diplomové práci.

% Review
\subsubsection{FP-16 Stav energie}

% Review
\begin{tabular}{lp{0.7\linewidth}}
    \textbf{Priorita:} & Should have \\
\end{tabular}

% Review
Aplikace by měla uživateli zobrazovat jeho stav energie, který je reprezentovaný body energie.
Učením se budou body energie snižovat. Pokud uživatel všechny body energie spotřebuje, nebude mu umožněno se dále učit. Body energie se budou v průběhu času obnovovat. Uživatelé s premium účtem budou mít neomezené množství energie, tedy učení nebude body energie snižovat.

% Review
Tento požadavek je inspirován analyzovanou aplikací Duolingo, která podobnou funkci nabízí. Tato funkce slouží primárně k motivaci uživatele, aby se učil pravidelně každý den a dále také k tomu, aby byl uživatel motivován si zakoupit premium účet.

% Review
\subsubsection{FP-17 Registrace a přihlášení pomocí Google a Apple}

% Review
\begin{tabular}{lp{0.7\linewidth}}
    \textbf{Priorita:} & Should have \\
\end{tabular}

% Review
Uživatel by měl mít možnost se registrovat a přihlásit do aplikace pomocí účtu Google nebo Apple. 

% Review
Tento požadavek vychází z analýzy konkurence, kde většina konkurenčních aplikací tuto možnost nabízí.

% Review
\subsubsection{FP-18 Zpřehlednění navigace v aplikaci}

% Review
\begin{tabular}{lp{0.7\linewidth}}
    \textbf{Priorita:} & Should have \\
\end{tabular}

% Review
Aplikace by měla poskytovat uživateli přehlednější a jednodušší způsob navigace mezi jednotlivými stránkami. Z výsledků uživatelského testování v mé bakalářské práci zabývající se stávající aplikací vyplynulo, že mají uživatelé problém s navigací mezi jednotlivými stránkami aplikace, zejména mezi jednotlivými úrovněmi mapy učení v daném kurzu. Dále, také z důvodu přidání nových stránek do aplikace, bude třeba zajistit, aby navigační prvky v aplikaci byly pro uživatel přehledné a snadné na používání.

% Review
Změnu navigace bude třeba provést zejména v hlavní navigační liště, bočním navigačním panelu a stránkou sloužící k přepínání mezi jednotlivými kurzy.

% Review
Cílem této změny je zpřehlednit aplikaci a zjednodušit uživateli navigaci mezi jednotlivými stránkami.

% Review
\subsubsection{FP-19 Tmavý režim}

% Review
\begin{tabular}{lp{0.7\linewidth}}
    \textbf{Priorita:} & Should have \\
\end{tabular}

% Review
Rozhraní aplikace by mělo být dostupné ve dvou režimech, ve světlém a tmavém. Uživatel by měl mít možnost v nastavení zvolit, zda chce používat světlý režim, tmavý režim, nebo režim dle systémového nastavení. Rozhraní aplikace by se pak mělo zobrazovat uživateli v tomto režimu.

% Review
\subsubsection{FP-20 Push Notifikace}

% Review
\begin{tabular}{lp{0.7\linewidth}}
    \textbf{Priorita:} & Should have \\
\end{tabular}

% Review
Aplikace by měla být schopna uživateli posílat push notifikace. % Write - zde přidat referenci na vysvětlení pojmu push notifikace
Tyto notifikace by měly uživatele upozorňovat v následujících situacích:

% Review
\begin{itemize}
    \item Pokud si uživatel nastavil každodenní připomenutí učení a ještě se v daný den neučil.
    \item Pokud se uživatel delší dobu v aplikaci nic neučil.
    \item Pokud uživateli končí předplatné aplikace a ještě ho neobnovil.
\end{itemize}

% Review
V rámci této práce se budu zabývat pouze zprovozněním zobrazování push notifikací uživateli ve frontendové části aplikace. Podrobným fungováním odesílání push notifikací se bude zabývat kolega Bc. Jakub Vondráček jeho diplomové práci, která se zabývá backendovou částí této aplikace. % Write - přidat referenci na backendovou diplomku

% Review
\subsubsection{FP-21 Připomínání učení}

% Review
\begin{tabular}{lp{0.7\linewidth}}
    \textbf{Priorita:} & Could have \\
\end{tabular}

% Review - tohle se ještě rozhodnout, jestli tam chci mít nebo ne
Uživatel by měl mít možnost si v aplikaci nastavit připomenutí, že se má v určitý čas učit. Aplikace by pak měla uživateli posílat push notifikace, které ho na učení upozorní, pokud se ještě v daný den neučil. Tato funkce je inspirována konkurenčními aplikacemi, které umožňují uživateli si nastavit každodenní připomenutí učení.

% Review
Dále by aplikace měla uživateli umožnit zapnout "chytré připomínání". Pokud uživatel tuto funkci zapne, nebude mu aplikace připomínat učení podle nastaveného času, ale na základě předchozího používání aplikace. Tato možnost by měla být v aplikaci nastavena jako výchozí.

% Review
\subsubsection{FP-22 Systém zobrazování reklam}

% Review
\begin{tabular}{lp{0.7\linewidth}}
    \textbf{Priorita:} & Should have \\
\end{tabular}

% Review
Ve stávající aplikaci se reklamy zobrazují pouze po úspěšném dokončení učení a to pouze tzv. bannerové reklamy. % Write - tady ověřit název reklam

% Review
Nově by aplikace měla umožňovat zobrazování dalších typů reklam, v závislosti na zařízení uživatele. Na mobilních zařízeních by aplikace měla primárně zobrazovat reklamy přes celou obrazovku. Na desktopových zařízeních by pak aplikace měla zobrazovat reklamy v podobě bannerů. Reklamy by měla aplikace zobrazovat v následujících případech.

\begin{itemize}
    \item Po skončení učení, ať už po úspěšném dokončení nebo po předčasném ukončení učení uživatelem.
    \item Pokud chce uživatel na účtu doplnit body energie.
    \item Pokud uživatel dostal odměnu s virtuální měnou a zvolil možnost dostání dvojnásobku odměny, pokud zhlédne reklamu.
\end{itemize}

% Review
Dále by aplikace měla uživateli v pravidelných intervalech zobrazovat reklamu na premium verzi aplikace.

% Review
Tento požadavek vychází z analýzy konkurenčních aplikaci, kde většina aplikací zobrazuje reklamy v průběhu používání aplikace a zobrazuje reklamy na premium verzi aplikace.

\subsubsection{FP-23 Nový systém předplatného}

% Review
\begin{tabular}{lp{0.7\linewidth}}
    \textbf{Priorita:} & Must have \\
\end{tabular}

% Review
Ve stávající aplikaci má uživatel možnost si zaplatit předplatné a tím získat přístup k pokročilejším funkcím. Jelikož v rámci této práce bude do aplikace přidána řada nových funkcí, je třeba tomu přizpůsobit i systém předplatného. Předplatné by nyní mělo být rozděleno do následujících kategorií.

% Review
\begin{itemize}
    \item \textbf{Uživatelé bez předplatného} budou mít přístup ke všem základním kurzům a kariérním cestám v aplikaci. Budou ovšem omezeni systémem energie, který jim umožní projít pouze omezený počet lekcí a cvičení za den. V aplikaci budou těmto uživatelům zobrazovány reklamy a budou mít omezený přístup k chatbotovi. Nebudou moci získat certifikáty o dokončení kurzů.
    \item \textbf{Uživatelé s předplatným} budou mít neomezený přístup ke všem kurzům a kariérním cestám aplikace. Nebudou jim zobrazovány reklamy, budou moci využívat chatbota s nižšími omezeními a budou moci získávat certifikáty o dokončení kurzů.
\end{itemize}

% Review
V rámci této funkce bude dále třeba zajistit lepší uživatelskou zkušenost a informovanost uživatelů, které funkce jsou dostupné v premium verzi aplikace. Pokud uživatel například chce využít funkci, která je dostupná pouze uživatelům s předplatným, měla by mu být zobrazena stránka, která ho bude informovat, že je daná funkce dostupná pouze v premium verzi aplikace.

\subsubsection{FP-23 Zkušební verze předplatného}

% Review
Aplikace by měla umožnit uživateli spustit zkušební verzi předplatného na určitou dobu zdarma. Tuto možnost by měli mít pouze uživatelé, kteří dříve zkušební verzi nikdy nevyužili.

% Review
\subsubsection{FP-24 Přizpůsobení avatara uživatele}

% Review
\begin{tabular}{lp{0.7\linewidth}}
    \textbf{Priorita:} & Should have \\
\end{tabular}

% Review
Uživatel by měl mít možnost si přizpůsobit svého avatara, neboli postavu na profilovém obrázku, pomocí položek, které zakoupil v obchodě za virtuální měnu. V rámci tohoto požadavku bude také třeba upravit vzhled profilové stránky uživatele, aby byl profilový obrázek viditelnější. Pro přizpůsobení avatara by aplikace měla uživateli poskytnout samostatnou stránku, kde může změny provádět. 

% Review - zde ještě rozhodnout, které položky chceme a které ne. Tehle seznam byl sepsaný podle toho, co má Duolingo, ale v něčem bychom se mohli třeba přiblížit tomu, jak to má třeba Habitica nebo Eatventure

Uživatel by měl mít možnost si u avatara přizpůsobit barvu pleti, velikost těla, obličej, barvu očí, vlasy, barvu vlasů, doplňky, barvu brýlí, vousy, barvu vousů, pokrývku hlavy, oblečení, barvu oblečení a barvu pozadí profilového obrázku.

% Review - zde ještě případně přepsat, podle toho, jestli budu dělat analýzu gamifikace
Tento požadavek vychází zejména z analýzy forem gamifikace a také analýzy konkurenční aplikace Duolingo, která podobnou funkci přizpůsobení si avatara uživatele nabízí.

% Review
\subsubsection{FP-25 Rozšíření obchodu a odemykání položek}

% Review
\begin{tabular}{lp{0.7\linewidth}}
    \textbf{Priorita:} & Could have \\
\end{tabular}

% Review
V aplikaci již existuje stránka obchodu, v němž si uživatel může zakoupit pomocí virtuální měny profilové obrázky.

% Review
Vzhledem k novým požadavkům aplikace je třeba tento obchod rozšířit tak, aby bylo možné kupovat i další položky, které jsou popsány v požadavku FP-16 Přizpůsobení avatara uživatele. % Write zde přidat odkaz na požadavek FP-16

% Review
Dále by v obchodu měla být zavedena nová funkce odemykání položek. Některé položky by měly být dostupné ke koupi pouze pokud uživatel dosáhl určitého ocenění v aplikaci. Do té doby by se položky měly uživateli zobrazovat jako zamčené.

% Review
U zamknutých položek by měl uživatel mít možnost si zobrazit informace o tom, co musí splnit, aby danou položku odemknul.

% Review
\subsubsection{FP-26 Uživatelský profil}

% Review
\begin{tabular}{lp{0.7\linewidth}}
    \textbf{Priorita:} & Should have \\
\end{tabular}

% Review
V aplikaci se již nachází stránka s uživatelským profilem, která obsahuje základní informace o daném uživateli. Po zavedení nových funkcí do aplikace je třeba tuto stránku aktualizovat tak, aby kromě již zobrazovaných informací o uživateli přehledně zobrazovala i následující informace.

% Review
\begin{itemize}
    \item Avatar uživatele, popsaný v požadavku FP-16 Přizpůsobení avatara uživatele. % Write - přidat odkaz na požadavek
    \item Liga, ve které se uživatel nachází. Popis lig je popsán v požadavku FP-19 Liga učení. % Write - přidat odkaz na požadavek
    \item Kurzy, které má uživatel zapsané a přehled, jaké úrovně v nich uživatel dosáhl.
\end{itemize}

% Review
Z důvodu přidání nových informací bude třeba změnit vzhled celé profilové stránky uživatele.

\subsubsection{FP-27 Zmrazení streaku}

% Review
\begin{tabular}{lp{0.7\linewidth}}
    \textbf{Priorita:} & Should have \\
\end{tabular}

% Review
Aplikace uživateli zobrazuje tzv. streak, neboli skóre, které reprezentuje počet dní v řadě, v nichž uživatel splnil alespoň jedno učení.

% Review
Uživatel by nově měl mít v obchodě s virtuální měnou možnost si zakoupit do zásoby několik položek tzv. zmrazení streaku.

% Review
Pokud uživatel bude mít zakoupenou jednu nebo více těchto položek, a nesplní daný den žádné cvičení, jeho streak se tím nezruší a jedna tato položka se spotřebuje. Díky tomu bude uživatel moci zachovat svůj streak i když některý den například zapomene se učit.

% Review
\subsubsection{FP-28 Oznámení o změně stavu streaku}

% Review
\begin{tabular}{lp{0.7\linewidth}}
    \textbf{Priorita:} & Should have \\
\end{tabular}

% Review
Pokud uživatel dokončí učení a jeho streak se zvýší na novou významnou hranici, například 100 dní, měla by mu aplikace zobrazit oznámení o dosažení této hranice.

% Review
Dále pokud uživatel ztratí streak, tedy je mu resetován na 0 dní, měla by mu tuto informaci aplikace oznámit po otevření aplikace. % Write - tohle líp přeformulovat

% Review
\subsubsection{FP-29 Liga učení}

% Review
\begin{tabular}{lp{0.7\linewidth}}
    \textbf{Priorita:} & Could have \\
\end{tabular}

% Review
V aplikaci by měly být zavedeny tzv. ligy učení, které by měly fungovat následujícím způsobem. Začátkem každého týdne by měl být každý uživatel zařazen do skupiny s omezeným počtem uživatelů, kteří se nachází ve stejné lize, tedy získali předchozí týden podobné množství bodů XP. Tato skupina je reprezentována žebříčkem uživatelů, který je seřazen podle počtu XP, které získali v momentálním týdnu.

% Review
Na konci každého týdne by měl být uživatel přesunut do vyšší ligy, pokud se nacházel v žebříčku uživatelů na prvních několika pozicích. Pokud se uživatel nacházel v žebříčku na posledních několika pozicích, měl by být přesunut do nižší ligy. Pokud se uživatel nacházel v žebříčku mezi těmito pozicemi, měl by zůstat v aktuální lize.

% Review
Tento požadavek vychází z analýzy konkurenčních řešení, zejména z analýzy aplikací Mimo a Duolingo, v nichž jsou podobné ligy implementovány.

% Review
\subsubsection{FP-30 Truhla s odměnou}

% Review
\begin{tabular}{lp{0.7\linewidth}}
    \textbf{Priorita:} & Should have \\
\end{tabular}

% Review
Po dokončení učení by uživateli měla být zobrazena stránka s truhlou, kterou bude moci uživatel otevřít. Po otevření truhly by měl uživatel získat náhodné množství virtuální měny, jako odměnu za jeho učení. Uživateli by následně mělo být nabídnuto, že může získat dvojnásobek této odměny, pokud zhlédne reklamu.

% Write - napsat, že tento požadavek vychází z analýzy gamifikace, pokud jsem jí prováděl - nebo jinak odůvodnit existenci tohoto požadavku

% Review
\subsubsection{FP-31 Denní výzvy}

% Review
\begin{tabular}{lp{0.7\linewidth}}
    \textbf{Priorita:} & Could have \\
\end{tabular}

% Review
Aplikace by měla uživateli každý den zobrazit několik denních výzev, které může splnit pro získání truhly s odměnou (vizte požadavek FP-19 Truhla s odměnou). % Write - přidat odkaz na požadavek
Tyto výzvy mají za cíl motivovat uživatele, aby se učil nad rámec jeho základního učení, aby vyzkoušel různé funkce aplikace, a aby si rozšířil svoje znalosti napříč různými kurzy a nejen v kurzech, které má zapsané.

% Review - zde ještě dopsat další, co mě napadnou
V aplikaci by měly být následující typy denních výzev:
\begin{itemize}
    \item Naučení se určitého počtu nových lekcí.
    \item Získání určitého počtu bodů XP.
    \item Učení se po určitý čas.
    \item Naučení se jedné lekce z určitého kurzu.
    \item Použití určité funkce aplikace, jako například zeptání se o radu chatbota nebo přizpůsobení avatara uživatele.
\end{itemize}

\subsubsection{FP-32 Certifikáty o dokončení kurzu}

% Review
\begin{tabular}{lp{0.7\linewidth}}
    \textbf{Priorita:} & Should have \\
\end{tabular}

% Review
Po dokončení kurzu by měl mít uživatel možnost si vygenerovat a stáhnout certifikát o dokončení kurzu. Certifikát by měl být ve formátu PDF a měl by obsahovat základní informace o dokončeném kurzu a uživateli, který kurz dokončil.

\subsubsection{FP-33 Sdílení pokroku mezi uživateli}

% Review
\begin{tabular}{lp{0.7\linewidth}}
    \textbf{Priorita:} & Should have \\
\end{tabular}

% Review
V aplikaci existuje základní funkce pro sdílení pokroku mezi uživateli. Tato funkce je ovšem příliš jednoduchá a s implementací nových funkcí je třeba ji přizpůsobit.

% Review
Aplikace by měla být přizpůsobena na sdílení následujících zpráv mezi uživateli.

% V rámci tohoto požadavku by taky měl být odstraněn infinite scroll. Místo toho by měla aplikace jednoduše posty za posledních 7 dní a další nezobrazovat.

% Review - zde vymyslet všechny další cases
\begin{itemize}
    \item Když uživatel dosáhne určitého streaku.
    \item Když uživatel skončí na předních příčkách v lize.
    \item Když uživatel dosáhne určité úrovně v kurzu.
    \item Když uživatel dosáhne určitého ocenění.
\end{itemize}

\subsection{Nefunkční požadavky}

\subsubsection{NP-1 Nový design aplikace}

% Review
\begin{tabular}{lp{0.7\linewidth}}
    \textbf{Priorita:} & Should have \\
\end{tabular}

% Review
Jelikož aplikace mění své zaměření z výuky jazyků na výuku softwarového inženýrství, je třeba změnit a přizpůsobit i samotný vzhled aplikace.

% Review
Měl by být vytvořen nový design uživatelského rozhraní, který bude aplikován jak na nové funkce, tak na funkce existující.

% Review
Základní komponenty používané při implementaci by dále měly být extrahovány do samostatné knihovny tak, aby bylo možné je využít i v dalších částech projektu, jako například na webových stránkách.

% Review
Tato knihovna komponentů by měla stavět na již existujících komponentech a měla by využívat již zavedené technologie, tedy knihovny React, TypeScript, Storybook a Material UI.


% Write - tohle možná přejmenovat na něco lepšího
\subsubsection{NP-2 Aktualizace kódu}

% Review
\begin{tabular}{lp{0.7\linewidth}}
    \textbf{Priorita:} & Should have \\
\end{tabular}

% Review
V kódu aplikace by měly být provedeny následující změny se záměrem zvýšit jeho kvalitu a udržitelnost a zjednodušit budoucí vývoj.

% Review
\begin{itemize}
    \item Aktualizace používaných knihoven, nástrojů a jejich konfigurací.
    \item Aktualizace používaných formatterů a linterů.
    \item Dockerizace frontendové části aplikace.
    \item Zavedení pravidel a konvencí pro nástroje, které využívají umělou inteligenci pro generaci kódu, dokumentačních komentářů a testů, s cílem podpoření vývoje aplikace a zvýšení efektivity práce.
    \item Odstranění existujícího mocku API, místo kterého by měla aplikace využívat testovací prostředí poskytované backendovou částí aplikace.
\end{itemize}

% Review
Tento požadavek byl vytvořen na základě výstupů analýzy existujícího řešení.

% Review
\subsubsection{NP-3 Zavedení automatizovaných testů}

% Review
\begin{tabular}{lp{0.7\linewidth}}
    \textbf{Priorita:} & Should have \\
\end{tabular}

% Review
Do frontendové části aplikace by měly být zavedeny automatizované testy s cílem zajištění správného fungování aplikace a zvýšení její kvality. Automatizované testy by se měly zaměřit na klíčové části aplikace, kterými jsou proces učení, registrace a přihlášení uživatele, výběr a zápis kurzů, a kupování předplatného.

% Review
Tento požadavek vychází jednak z analýzy existujícího řešení a jednak z návrhu úprav do budoucna v mé bakalářské práci zabývající se zmíněným řešením. % Write - přidat referenci na bakalářskou práci

% Review
\subsubsection{NP-4 Odstranění přebytečných funkcí}

% Review
\begin{tabular}{lp{0.7\linewidth}}
    \textbf{Priorita:} & Must have \\
\end{tabular}

% Review
Jelikož se aplikace bude zaměřovat na výuku softwarového inženýrství, bude třeba odstranit funkce týkající se původní domény výuky jazyků. Mezi tyto funkce patří následující.

% Write - tento seznam extrahovat do kapitoly analýza stávajícího systému a zde na to pouze odkázat.
\begin{itemize}
    \item Vyhledávání slovíček
    \item Přidávání slovíček do seznamu oblíbených slovíček
    \item Cvičení zaměřená na výuku jazyků a která nejsou využitelná pro výuku softwarového inženýrství
    \item Rozdělení typů lekcí dle aspektu výuky jazyků, například gramatické lekce, lekce slovní zásoby apod.
    \item Stránka slovíčka
    \item Výběr preferovaných témat kurzu. Tato funkce v původní aplikaci sloužila jako způsob přeskočení skupin lekcí slovíček, které se netýkaly zájmů a cílů uživatele. Uživatel tak například mohl vypnutím tématu "Příroda a životní prostřední" jednoduše ze svého učení vyřadit všechny lekce týkající se tohoto tématu. V nové aplikaci ovšem bude struktura kurzů a domény změněna natolik, že tato funkce již nebude relevantní a její ponechání by mohlo mít za následek značné navýšení komplexity celého systému.
    \item Přidávání lekcí do seznamu oblíbených lekcí. Žádná z analyzovaných aplikací tuto funkci nenabízí. Funkce v rámci původní domény výuky jazyků mohla být užitečná, nicméně nyní by zbytečně mohla komplikovat uživatelské rozhraní.
    \item Odebrání stránky lekce. Tato stránka dříve zobrazovala například seznam vyučovaných slovíček v dané lekci. Zároveň bylo možné na této stránce přidávat lekci do seznamu oblíbených lekcí. V rámci domény výuky softwarového inženýrství žádná z analyzovaných aplikací samostatnou stránku pro lekci neobsahuje. Nyní ovšem stránka žádné výhody nemá a pouze komplikuje uživatelské rozhraní.
    \item Vytváření vlastních lekcí slovíček. Původní aplikace umožňovala uživatelům seskupovat slovíčka v aplikaci do vlastních lekcí. Tato funkce již není relevantní z důvodu nové domény a struktury kurzů.
\end{itemize}

% Review
\subsubsection{NP-5 Automaticky generované třídy API}

% Review
\begin{tabular}{lp{0.7\linewidth}}
    \textbf{Priorita:} & Should have \\
\end{tabular}

% Review
V rámci kódu frontendu by měly být automaticky generovány třídy reprezentující backendové API na základě openapi % Write - zjistit jak se píše správně open api specifikace
specifikace poskytované backendovou částí aplikace.

% Review
Požadavek vychází z analýzy existujícího řešení. Jeho cílem je zjednodušení implementace a zajištění konzistence mezi frontendovou a backendovou částí aplikace.

\subsubsection{NP-6 Animace a mikro interakce}

% Review
\begin{tabular}{lp{0.7\linewidth}}
    \textbf{Priorita:} & Should have \\
\end{tabular}

% Review
Do aplikace by měly být přidány základní animace s cílem zlepšit uživatelskou zkušenost a udržet pozornost uživatelů. Uživatelské rozhraní by tak mělo působit vizuálně poutavějším dojmem a mělo by výrazněji reagovat na interakce s uživatelem.

% Review
Tento požadavek vychází zejména z analýzy konkurenčních řešení, ve které konkurenční aplikace jako Mimo nebo Scrimba využívají animace a efekty.

\subsubsection{NP-7 Sestavování a nasazování PWA}

% Review
\begin{tabular}{lp{0.7\linewidth}}
    \textbf{Priorita:} & Should have \\
\end{tabular}

% Review
Aplikace by měla být přizpůsobena tak, aby bylo možné ji sestavit a nasadit jako PWA, neboli progresivní webovou aplikaci. Tomu by měly být přizpůsobeny i devops procesy, v rámci kterých by měla být aplikace automaticky sestavována a nasazována jako PWA.

%\subsubsection{(NP-8 Sběr statistik o uživatelském chování)}
% Review - tenhle požadavek spíš dělat nebudu.. zabere to zbytečně moc času.. RESP můžu dělat třeba základní sběr statistik o tom, jak dlouho zaberou uživatelům jednotlivá cvičení, ale to bych dal rovnou do požadavku o změně učení.. DOBRÝ BY ALE BYLO, ŽE BYCH DO KAPITOLY TESTOVÁNÍ MOHL ZAPSAT I KVANTITATIVNÍ STUDII - nicméně na to bych stejně potřeboval aby to vyplnilo hodně uživatelů, takže bych stejně asi neměl žádné výsledky té studie..